\documentclass[../main.tex]{subfiles}
\begin{document}

\chapter{Cross-Week Concept Threads}
\label{ch:crossweek}

The course is taught week by week, but it is built as two threads that each run the length of the programme. This chapter follows both threads end to end, showing how the output of one week becomes the input of the next. It is the chapter to read when a question spans weeks rather than sitting inside one.

\section{The Conceptual Track}
\label{sec:thread-conceptual}

The conceptual thread begins in Week~1 with the language of research and the shape of its reasoning. The vocabulary fixed there, the hypothesis, the variable, and the three study types (\autoref{sec:language-research}), is what later weeks manipulate, and the hourglass logic of deduction down to measurement and induction back to generalisation (\autoref{sec:hourglass}) is the inferential frame the whole track inherits. Nothing here is yet a plan, but everything the plan will be made of is named, which is why the track starts with definitions rather than activities.

Week~2 turns that vocabulary into a design. The hypothesis becomes a research objective, the objective generates research questions, and the questions are sequenced by the research framework (\autoref{sec:objectives-rqs}). The decisive step is operationalisation (\autoref{sec:operationalization}), where the abstract concepts of Week~1 are pinned to exact definitions, stated assumptions, and named variables, because this is the point at which the conceptual design hands concrete quantities to the technical design (\autoref{sec:strategy-materials}). The split between the conceptual question of what and why and the technical question of how and when (\autoref{sec:conceptual-technical}) is the organising move of the week, and it only works because the language of Week~1 made the concepts precise enough to operationalise.

Week~6 closes the thread by writing it down. The operationalised plan becomes the body of the Research Methodologies report, whose sections are the research questions and the planning built earlier (\autoref{sec:reports}), and the literature that seeded the objective in Week~1 (\autoref{sec:research-ideas}) returns as the reference list that must meet the sourcing and style rules of Week~6 (\autoref{sec:referencing}). Read in sequence, the conceptual track is a single sentence stretched across the course: a precise question, made into a designed plan, then communicated to a reader who can act on it.

\section{The Empirical Track}
\label{sec:thread-empirical}

The empirical thread begins in Week~1 with the six steps of a research project (\autoref{sec:research-process}), and two of those steps carry the whole track. The sample step asks which portion of the problem will be studied, and the measure step asks how observations will be recorded, and both are stated only in outline at this stage. The track is the progressive sharpening of these two steps into rigorous, checkable, and managed practice over the following weeks.

Week~3 makes the sample and measure steps exact. The sample step becomes the population-to-sample funnel and the statistics of representativeness (\autoref{sec:sampling}), with the probability boundary deciding what can be generalised and the standard error quantifying precision (\autoref{sec:sample-types}). The measure step becomes the treatment of data sources, the legal categories that constrain them, the qualitative and quantitative split, and the quality checks and logbook that protect them (\autoref{sec:data-sources}). What was a one-word step in Week~1 is now a method with rules and a record.

Week~4 then interrogates the data the empirical track has produced. Verification asks whether the method was applied correctly and validation asks whether the result is true (\autoref{sec:vv}), and sustainable research adds a third interrogation, weighing the value, reproducibility, and resource cost of the work through the Engineering for One Planet framework (\autoref{sec:eop}). The resource axis in particular reaches back to the standard error of Week~3, since it demands that the sample be no larger than the objective justifies. Week~5 finally manages both the data and the project, as the logbook of Week~3 grows into the six-section data management plan (\autoref{sec:dmp}) and the planning outlined in Week~2 grows into a Work Breakdown Structure and a Gantt chart with a critical path (\autoref{sec:project-mgmt}). The empirical track is therefore a tightening loop, from a rough step, to a rigorous method, to a checked result, to a managed system, with each week formalising what the previous one left informal.

\section{Where the Tracks Meet}
\label{sec:thread-meet}

The two tracks are not parallel lines that never touch. Sustainable research sits on the empirical side yet feeds project management and data management, so the value and resource questions of one track shape the planning decisions of the other (\autoref{sec:eop}). Data management collects the output of data, sustainable research, and technical design at once (\autoref{sec:dmp}), which is why the course map draws it as the node where the red and green chains converge. The Research Methodologies report is the final meeting point, since its body carries the conceptual plan while its conduct, sourcing, and data handling carry the empirical discipline, and the report guide that follows (\autoref{ch:reportguide}) shows section by section how both tracks are written into one document.

\end{document}
